\documentclass[12pt]{article}

\usepackage{fontawesome}
\usepackage{hyperref}
\usepackage{xurl}
\usepackage[tagged, highstructure]{accessibility}
\usepackage{bookmark}
\usepackage{enumitem}

\hypersetup{
    colorlinks=false,
    pdfborder={0 0 0},
    pdftitle={Common commands,logging into a server and sharing files},
    pdfauthor={Adrianna Holden-Gouveia},
    pdfsubject={Introduction to Linux - Lab Assignment},
    pdflang={en-US},
    pdfkeywords={lab assignment, Common commands,logging into a server and sharing files}
}

% Configure PDF bookmarks for navigation
\bookmarksetup{
    numbered,
    open,
}

% Configure list spacing for better accessibility
\setlist{nosep}



\title{Common commands,logging into a server and sharing files}
\author{
        Adrianna Holden-Gouveia \\
        Website: \href{https://aholdengouveia.name}{aholdengouveia.name}\\ 
        \date{\vspace{-5ex}}
        %Email: \href{mailto:admin@aholdengouveia.name}{admin@aholdengouveia.name} \\
%        \faLinkedin{: aholdengouveia} \\
        \faGithub {: aholdengouveia} \\
%        \faTwitter {: aholdengouveia} \\
        }

%S\date{\today}


\begin{document}    

\maketitle

\section*{Accessibility Notice}
This document is also available in HTML format at:

Link: \href{https://aholdengouveia.name/IntroLinux/labs/commoncommands.html}{\textbf{https://aholdengouveia.name/IntroLinux/labs/commoncommands.html}}

The HTML version provides enhanced accessibility features including keyboard navigation, screen reader support, responsive design, dark mode support, and high contrast options.


%\begin{abstract}
%This is a template for Linux Administration Lab work
%\end{abstract}
%\tableofcontents

\section*{Objectives:}
\begin{itemize}
    \item Demonstrate logging on, moving files to and from the local system to the server
\end{itemize}
\section*{Complete the following problems}

References, a video, a PowerPoint and some notes are available at my website
Link: \href{https://www.aholdengouveia.name/IntroLinux/commoncommands.html}{\textbf{https://www.aholdengouveia.name/IntroLinux/commoncommands.html}}

The lab report will include screen shots as evidence of completion.  

If you are my current student, Resources for the server can be found from the menu item "Server Access Information" on your LMS



\subsection*{Steps for the lab}
    \begin{enumerate}
        \item Use puTTY or SSH to connect to the server (or do this on you own Linux system)  If you need help using SSH or SCP I made a video on how to do that  on my YouTube channel.  SCP?  SSH-ut the front door! Link: \href{https://youtu.be/w3yOQFZJMTo}{\textbf{https://youtu.be/w3yOQFZJMTo}}
        \item Login to your account (Information from Instructor posted to LMS)
        \item Test by  logging off and logging in again
        \item Use ls and ls -a to see what is in your home directory. Make three empty files called foo1 foo2 and foo3 using touch and use ls again. Take a screen shot.
        \item Create a sub-directory called Lab0 and cd to the subdirectory and create three more files using touch called bar1 bar2 and bar3.  Use ls and take a screen shot
        \item Create three text files in the local system using notepad, call them localfoo1, localfoo2 and localfoo3.  Put different texts in each file. Add your name and date to the top of each file
        \item Using WinSCP, FileZilla or other SCP client of your choice in the local system move the three localfoo files to the server. Take a screen shot.
        \item Use ls on the server to show that the files are in the server; use cat to display one of the files and take a screen shot.
    \end{enumerate}



\section*{Deliverables}
One file should be submitted, it should have the lab title, your name and date, and include each of the screen shots with a one sentence description of what is in the screen shot.
\end{document}